\documentclass[10pt]{article}
\usepackage[a4paper,margin=18mm]{geometry}
\usepackage{amsmath,amssymb,booktabs,array,enumitem}
\usepackage[T1]{fontenc}\usepackage{lmodern}
\usepackage{tikz}\usetikzlibrary{shapes.gates.logic.US,positioning,arrows.meta}
\usepackage[hidelinks]{hyperref}
\setlength{\parindent}{0pt}\setlength{\parskip}{4pt}
\setlength{\emergencystretch}{2em}
\newcommand{\VDD}{V_{DD}}
\newcommand{\mos}[4]{\begin{scope}[shift={(#1,#2)}]
\draw (0,.24)--(0,-.24);\draw (0,.20)--(.25,.20)--(.25,.5);
\draw (0,-.20)--(.25,-.20)--(.25,-.5);\draw (-.16,.23)--(-.16,-.23);
\ifnum#4=1\draw (-.26,0) circle (.06);\draw (-.75,0)--(-.32,0);\else\draw (-.75,0)--(-.16,0);\fi
\node[anchor=east,font=\scriptsize] at (-.76,0) {$#3$};\end{scope}}
\newcommand{\ground}[2]{\draw (#1,#2)--(#1,{#2-.12});\draw ({#1-.17},{#2-.12})--({#1+.17},{#2-.12});\draw ({#1-.11},{#2-.20})--({#1+.11},{#2-.20});}
\newcommand{\rail}[3]{\draw (#1,#2)--(#1,{#2+.14});\node[above,font=\scriptsize] at (#1,{#2+.14}) {$#3$};}

\usetikzlibrary{calc,patterns}
\newcommand{\hm}[4]{\begin{scope}[shift={(#1,#2)}]
\draw(-.24,0)--(.24,0);\draw(-.2,0)--(-.2,-.25)--(-.5,-.25);
\draw(.2,0)--(.2,-.25)--(.5,-.25);\draw(-.23,.16)--(.23,.16);
\ifnum#4=1\draw(0,.26)circle(.06);\draw(0,.32)--(0,.75);\else\draw(0,.16)--(0,.75);\fi
\node[font=\scriptsize,above right]at(.30,.05){$#3$};\end{scope}}
\newcommand{\tg}[3]{\draw(#1-.24,#2-.13)rectangle(#1+.24,#2+.13);\node[above,font=\tiny,inner sep=2pt]at(#1,#2+.13){$e=#3$};}
\begin{document}
\begin{center}{\large\bfseries EE3032 Introduction to VLSI -- Midterm}\end{center}
\textbf{1. (35 points)} Answer the following questions briefly.
\begin{enumerate}[label=(\alph*),leftmargin=7mm,itemsep=2pt,topsep=2pt]
\item List the process for transferring the mask pattern to the surface of a silicon region. (3 points)
\item List two approaches for alleviating the crosstalk effect of two parallel interconnects. (2 points)
\item What silicon on insulator process is and explain how it can reduce the leakage current? (4 points)
\item List two methods for the oxidation in silicon. (2 points)
\item List two features of static CMOS circuit. (2 points)
\item List two components of CMOS circuit, which result in static power. (2 points)
\item In a logic gate, we usually implement the pull-up network with PMOS transistors and the pull-down network with NMOS transistors. Why? (2 points)
\item List two through-silicon via (TSV) process technologies. (2 points)
\item In a CMOS gate, the transistors connected to the gate output suffer from the severest body effect. Why? (2 points)
\item What is antenna effect? List two design methods for the fixing of antenna effect violation. (6 points)
\item List three bonding technologies for 3D ICs. (3 points)
\item The output capacitance of a CMOS logic gate in a CMOS circuit is comprised of three components. List the three components. (3 points)
\item List two types of FinFET structures. (2 points)
\end{enumerate}
\textbf{2. (15 points)} (a) Design an 8:1 multiplexers using transmission gate. (10 points) (b) Design a 3-input NAND gate using an 8:1 multiplexer. (5 points)

\textbf{3. (15 points)} Find CMOS realizations of the truth table shown in Table 1, where ABCD are the inputs and $F$ is the output.
\begin{center}\begin{tabular}{ccccc}\toprule $A$&$B$&$C$&$D$&$F$\\\midrule
0 & 0 & 0 & 0 & 0\\
0 & 0 & 0 & 1 & 0\\
0 & 0 & 1 & 0 & 0\\
0 & 0 & 1 & 1 & 1\\
0 & 1 & 0 & 0 & 0\\
0 & 1 & 0 & 1 & 1\\
0 & 1 & 1 & 0 & 0\\
0 & 1 & 1 & 1 & 1\\
1 & 0 & 0 & 0 & 0\\
1 & 0 & 0 & 1 & 1\\
1 & 0 & 1 & 0 & 1\\
1 & 0 & 1 & 1 & 1\\
1 & 1 & 0 & 0 & 0\\
1 & 1 & 0 & 1 & 1\\
1 & 1 & 1 & 0 & 1\\
1 & 1 & 1 & 1 & 1\\\bottomrule\end{tabular}\\Table 1: Truth table.\end{center}\vfill\small Archive attribution: 2021. The photographs do not print a date. Editable reconstruction of the complete two-page paper; original wording and point labels are retained.
\newpage
\begin{center}\begin{tikzpicture}[x=.85cm,y=.85cm,font=\scriptsize]
\draw[dashed](.15,4.5)rectangle(1.65,5.4);\node[align=center]at(.9,4.95){P\\network};\draw(.9,5.4)--(.9,5.8);\node[above]at(.9,5.8){1.2 V};
\draw(.9,4.5)--(.9,4.1)--(1.25,4.1)--(1.25,3.5);\mos{1}{3}{N3}{0}
\draw(1.25,2.5)--(1.25,2);\mos{1}{1.5}{N1}{0};\ground{1.25}{1}
\draw(-.75,2.25)--(.05,2.25)--(.05,3)--(.25,3);\draw(.05,2.25)--(.05,1.5)--(.25,1.5);\fill(-.75,2.25)circle(.045);\node[left]at(-.75,2.25){$V_{in}$};
\draw(1.25,4.1)--(4.6,4.1);\fill(4.6,4.1)circle(.045);\node[right]at(4.6,4.1){$V_{out}$};
\hm{3.25}{2.25}{N2}{0};\draw(1.25,2.25)--(2.75,2);\fill(1.25,2.25)circle(.035);\draw(3.75,2)--(4.1,2)--(4.1,2.6);\node[above]at(4.1,2.6){1.2 V};\draw(3.25,3)--(3.25,4.1);\fill(3.25,4.1)circle(.03);\node at(2.2,.25){(a)};
\begin{scope}[shift={(7.2,1.6)},x=.85cm,y=1.7cm]
\draw[->](0,0)--(4.25,0)node[right]{$t$};\draw[->](0,0)--(0,1.65)node[above]{Voltage};
\draw[thick](0,1.2)--(1.65,1.2)--(2.85,0)--(4,0);\node at(.75,1.36){$V_{out}$};
\draw[thick](1.1,0)--(2.55,1.2)--(4,1.2);\node at(3.7,1.4){$V_{in}$};
\draw[dashed](0,.4)--(1.58,.4)--(1.58,0);\draw[dashed](0,.8)--(2.07,.8)--(2.07,0);\node[left]at(0,.4){0.4 V};\node[left]at(0,1.2){1.2 V};\node[below]at(1.58,0){$t_1$};\node[below]at(2.07,0){$t_2$};\node at(2.07,.9){$x$};\end{scope}
\node at(8.75,.25){(b)};
\end{tikzpicture}\\Figure 1: (a) A CMOS logic gate (b) waveforms of $V_{in}$ and $V_{out}$.\end{center}
\textbf{4. (15 points)} Fig. 1(a) shows a CMOS logic gate. Assume that the threshold voltage of the NMOS transistor is 0.4 V; the turn on resistance of N1, N2, and N3 is $R$, $2R$, and $R$, respectively.
\begin{enumerate}[label=(\alph*),leftmargin=7mm]
\item If the initial voltage of the output $V_{out}$ is 1.2 V and an input stimulus $(0\to1.2\text{ V})$ as shown in Fig. 1(b) is applied, how much voltage of the input stimulus ($V_{in}$) at $x$ is? (9 points)
\item Identify the operation region of the transistors N1, N2, and N3 from $t_1$ to $t_2$. (6 points)
\end{enumerate}
\textbf{5. (10 points)} What function does the circuit shown in Fig. 2 realize? Does it have limitations?
\begin{center}\begin{tikzpicture}[x=.9cm,y=.9cm,font=\small]
\mos{3}{3}{\phantom{B}}{1}\mos{3}{1.5}{\phantom{B}}{0}
\draw(3.25,3.5)--(3.25,4)--(.15,4)node[left]{$A$};\draw(3.25,1)--(3.25,.5)--(.15,.5)node[left]{$\bar A$};
\draw(3.25,2.5)--(3.25,2)--(4,2)node[right]{$y$};\draw(3.25,2)--(3.25,2);\draw(3.25,2)--(3.25,2);\draw(3.25,2)--(3.25,2);\draw(3.25,2)--(3.25,2);
\draw(.15,2)--(1.28,2);\node[left]at(.15,2){$B$};\draw(1.28,1.76)rectangle(1.78,2.24);\draw(1.78,2)--(3.25,2);
\draw(1.53,2.24)--(1.53,2.65);\draw(1.53,2.76)circle(.055);\draw(1.53,2.82)--(1.53,4);\draw(1.53,1.76)--(1.53,.5);
\draw(.65,2)--(.65,1.25)--(2.1,1.25)--(2.1,3)--(2.25,3);\draw(2.1,1.5)--(2.25,1.5);\fill(.65,2)circle(.03);\fill(2.1,1.5)circle(.03);\fill(3.25,2)circle(.03);
\end{tikzpicture}\\Figure 2: A CMOS circuit.\end{center}
\textbf{6. (10 points)} Complete the following questions.
\begin{enumerate}[label=(\alph*),leftmargin=7mm]
\item Implement the following Boolean function: $F=X(Y+Z)+XW$ using CMOS logic. Draw the schematic diagram. (5 points)
\item We have a CMOS logic gate realizing $F=\overline{X(Y+Z)}$ and the sizes of all the NMOS and PMOS transistors of the CMOS logic gate are $W/L$. Design an equivalent inverter for analyzing the rise time ($t_r$) and fall time ($t_f$) of the logic gate with input transistors $XYZ:101\to100\to101$. (5 points)
\end{enumerate}
\vfill\small The six printed question totals sum to 100 points. Figures are redrawn without changing connectivity; Fig. 1(b) is qualitative.
\end{document}
